\subsection{Two obstructions to strict integral replacement constructions}

These statements concern proposed constructions for C.3; neither is a
counterexample to the common-source or terminal-groupoid assertion.
Put $S=\Gamma_{\Z}(T)$ and let $F$ be the standard theory. For $N>1$,
write $A_N=S+F^{\mathrm{weights}\geq N}$ for its standard tail
subtheory. It is admissible, as in Statement~\ref{sa08:C13}, and
$A_N\to F$ is an actual morphism.

In the next statement a \emph{free linear-disk PD DGA} over $S$
means the free strictly graded commutative PD algebra on a direct
sum of cochain disks
\[
 D^q(w)=(\Z x^q\xrightarrow{1}\Z y^{q+1}),
\]
with specified nonnegative weights. Odd generators are exterior;
even generators have divided powers. This is a restriction on
the differential: it does not mean an arbitrary semifree DGA
whose generator differentials may involve earlier generators
or decomposable expressions.

\begin{proposition}[C14: \status{Proved}; free linear-disk replacements fail]
\label{sa08:C14}
No simplicial object $P$ whose levels are free linear-disk PD
DGAs, and which satisfies eventual vertical exactness, admits
a degreewise surjective morphism $P\to A_N$.
\end{proposition}
\begin{proof}
Suppose a level contains an odd-to-even disk, with $q$ odd.
Setting all other generators to zero retracts that level onto
\[
 D=S\otimes\Lambda(x)\otimes\Gamma(y),\qquad dx=y,\quad dy=0.
\]
For every $m\geq2$,
\[
 d\bigl(x\gamma_{m-1}(y)\bigr)=m\gamma_m(y).
\]
Consequently $H^{m(q+1)}(D)$ contains
$(S/mS)\gamma_m(y)$. Already $m=2$ gives nonzero classes
$T^{[r]}\gamma_2(y)$ for all $r\geq0$, in a fixed cochain
degree and unbounded weights. They survive under the
retraction. This contradicts eventual vertical exactness,
even in this one simplicial level.

Only even-to-odd disks can therefore occur. Every degree-one
element of a free algebra on these disks is a sum of terms
$a\,dx$, with $a$ of cochain degree zero and $x$ a degree-zero
generator. Consider its image in weight $N$ in $A_N$.
If the weight of $x$ is less than $N$, its image is a scalar
coefficient, so $dx$ maps to zero. If its weight equals $N$,
the coefficient $a$ has weight zero and maps to a scalar,
so the image of $a\,dx$ is closed. Larger weights cannot
occur. This explicitly includes products of degree-zero
and degree-one generators; their differentials vanish in
this weight in this target.

In simplicial degree two, use cumulative coordinates
$X,Y$ for
$F_2=S\{X,Y\}\otimes\Lambda(dX,dY)$. Then
$X^{[N-1]}dY\in A_N^{N,1}$ and
\[
 d\bigl(X^{[N-1]}dY\bigr)
 =X^{[N-2]}dX\,dY\ne0.
\]
It cannot be in the indicated image, contradicting
degreewise surjectivity.
\end{proof}

Thus the canonical free/cotriple simplicial construction
in the algebraic theory of strict PD DGAs cannot provide
an admissible replacement of $A_N$. A horizontal trivial
fibration is degreewise surjective, so the same applies
to any replacement of this particular free form.
Statement~\ref{sa08:C14} does not address retracts of free
objects, semifree algebras with nonlinear differentials,
or common-source maps that are not trivial fibrations.

\begin{proposition}[C15: \status{Proved}; free coefficient-tail resolutions fail]
\label{sa08:C15}
Let $I_N\subset S$ be the full weight tail at least $N>1$,
and let $Q\to R=S/I_N$ be a graded cellular free simplicial
PD $S$-algebra resolution. Here cellular means obtained
from $S$ by successively adjoining free PD simplex
generators along their simplicial boundaries; the
augmentation is a weak equivalence in every weight.
For the special PD differential module,
defined by $d\gamma_m(a)=\gamma_{m-1}(a)\,da$, one has
\[
 \pi_1\bigl(\Omega^1_{\mathrm{PD}}(Q/S)\bigr)^p\ne0
 \qquad\text{for every prime }p>N.
\]
Therefore $\mathrm{DR}_{\mathrm{PD}}(Q/S)$ fails eventual
horizontal acyclicity in form degree one.
\end{proposition}
\begin{proof}
Fix a prime $p>N$. Define the subgroup $J_p\subset S$
by killing every $T^{[n]}$ with $n\geq N$, $n\ne p$,
and killing $pT^{[p]}$. It is an ideal: the only
product requiring attention lands in weight $p$ and
has coefficient $\binom pa$, divisible by $p$ for
$0<a<p$.

It is also a PD ideal. A term of weight $p$ in a
higher divided power of an element supported in
weights at least $N$ has coefficient
\[
 \frac{p!}
 {\prod_d(d!)^{m_d}\prod_d m_d!},
 \qquad
 \sum_d d\,m_d=p,\quad \sum_d m_d\geq2,\quad d\geq N.
\]
Every denominator factorial has argument strictly
less than $p$, so the coefficient is divisible by
$p$. Terms in other tail weights are already killed.
A term involving $pT^{[p]}$ cannot contribute to
weight $p$ in a higher divided power. The PD addition
and scalar formulas finish the check.

Hence the strict PD extension
\[
 E_p=S/J_p\longrightarrow R
\]
has kernel $M_p=\mathbb F_pT^{[p]}$. This kernel has
square zero and zero higher divided powers, and
positive-weight elements of $R$ act on it by zero.
The extension has no $S$-algebra section: the
coefficient $T^{[p]}$ is zero in $R$ but nonzero
in $E_p$.

Put $\Omega=\Omega^1_{\mathrm{PD}}(Q/S)$ and
$L=R\otimes_Q\Omega$. We first prove that
$\Omega\to L$ is a weak equivalence. Filter $\Omega$
by the stages at which the free generators of $Q$
were adjoined, extending coefficients to the final
algebra $Q$. A simplex generator of dimension $n$
adds one differential-module cell. Its boundary is
the differential of the prescribed algebra boundary,
and its quotient module is
\[
 Q\otimes_{\Z}
 \bigl(\Z[\Delta^n]/\Z[\partial\Delta^n]\bigr),
\]
with its generator weight. These filtration
inclusions are split in each simplicial degree as
$Q$-modules. The corresponding quotient for $L$
has $R$ in place of $Q$. Normalization and
Eilenberg--Zilber identify its comparison with
tensoring $NQ\to NR$ by a bounded free cell complex.
It is a weightwise quasi-isomorphism. Induction over
successor stages and exactness of filtered colimits
over limit stages prove the claim. This uses the
cellular presentation, not just degreewise freeness.

The module $H_0(NL)$ represents special PD
derivations $R\to M$ over $S$, by the coequalizer
defining $\pi_0Q=R$. All such derivations vanish
because $S\to R$ is surjective. Thus $H_0(NL)=0$.

Freeness gives an $S$-PD lift
$\alpha:Q_0\to E_p$ of the augmentation. The
difference
\[
 \delta=\alpha d_0-\alpha d_1:Q_1\longrightarrow M_p
\]
is a special PD derivation. Indeed products use
$M_p^2=0$, and divided powers use
\[
 \gamma_m(e+u)-\gamma_m(e)
 =\gamma_{m-1}(e)u\qquad(u\in M_p),
\]
since all higher divided powers of $u$ vanish.
It is therefore represented by an $R$-linear map
$L_1\to M_p$. It is normalized, and its cochain
coboundary vanishes by the simplicial identities.
If $\delta=h d_0-h d_1$ for a special derivation
$h:Q_0\to M_p$, then $\alpha-h$ is again an
$S$-PD algebra map and coequalizes the two faces.
It descends to the impossible $S$-PD section
$R\to E_p$. Consequently
\[
 H^1\operatorname{Hom}_{R,\mathrm{graded}}(NL,M_p)
 \ne0.
\]

This forces $H_1(NL)^p\ne0$. Otherwise $\delta$
kills the whole kernel of
$d:(NL)_1^p\to(NL)_0^p$. Since $H_0(NL)=0$,
the differential is surjective in every weight.
The map $\delta$ then factors through
$h:(NL)_0^p\to M_p^p$. Extend $h$ by zero in
other weights. It is $R$-linear: lift any
lower-weight element of $(NL)_0$ through $d$,
and use the $R$-linearity of $\delta$ and the
zero positive-weight action on $M_p$.
This makes $\delta$ a coboundary, a contradiction.
The weak equivalence $\Omega\to L$ proves the
assertion.
\end{proof}

The proof deliberately uses special square-zero
extensions with zero higher PD operations on their
kernels. General Beck modules of PD algebras can
carry an additional Frobenius operation; their
cotangent complex is not silently identified with
the special differential module used here
\cite[commutative-operad example]{BDokasFranklandIkonicoff2024}.
The ruled-out construction is precisely a free
horizontal resolution of the coefficient tail
followed by relative PD de Rham. Nonfree
replacements with further relations or higher
cochain generators remain outside its scope.

For a finite exact arithmetic certificate of the
$J_p$ calculation,
\texttt{B\_check\_prime\_pd\_ideal.py} checks
$2\leq N\leq\min(8,p-1)$ for all primes $p\leq43$:
79 pairs, 1,520 binomial products, and 35,431
PD partition coefficients. Every coefficient
passes the stated divisibility test.
The full receipt is
\texttt{B\_prime\_pd\_ideal\_checks.json}.
